\begin{figure}[tbp]
\centering
\begin{tikzpicture}[>=Latex,scale=.95,every node/.style={font=\small}]
  \draw[densely dashed,gray] (0,0) ellipse (4.4 and 2.05);
  \node[fill=white,inner sep=2pt] at (3.7,1.2) {$S_R$};
  \filldraw[fill=black!13,draw=black,thick] (-1.75,0) ellipse (1.12 and 1.08);
  \draw[densely dashed] (-1.92,-.02) ellipse (.48 and .42);
  \node at (-1.92,-.02) {$S$};
  \node at (-1.75,.72) {$\cD_{c_b}$};
  \filldraw[fill=white,draw=black,thick] (1.5,.65) ellipse (.78 and .53);
  \node at (1.5,.65) {$\cW_{c_w}$};
  \draw[->,thick] (-.63,0) -- (.18,0) node[midway,above] {$\nu_b$};
  \draw[->,thick] (1.5,.12) -- (1.5,-.65)
        node[midway,right] {$\nu_w$};
  \node at (.15,-1.05) {$\Omega_R$};
  \node[fill=white,inner sep=3pt] at (-1.75,-1.42) {$H+P_B=c_b$};
  \node at (1.5,1.55) {$H-P_B=c_w$};
  \draw[->] (4.4,0) -- (5.2,0);
  \node[right] at (5.2,0) {end \(e\)};
\end{tikzpicture}
\caption{A schematic of the auxiliary exterior; one black and one
possible white frontier component are shown. The white family may be
empty. The black region encloses the original boundary.
Both normals point into the retained exterior, but its black and white
faces have prescribed expansion for \(K\) and \(-K\), respectively.
The drawing represents incidence and orientations, not the topology
or dimension of the hypersurfaces.}
\label{four:fig:two-frontiers}
\end{figure}
